\chapterimage{chapter_head_1.pdf}
\chapter{Introduction to Computer Architecture}

This chapter covers the fundamental concepts of Computer Design and Architecture (CDA) as presented in the textbook and the introductory course materials.

\section{What is Computer Design and Architecture (CDA)?}

\begin{definition}[Computer Design and Architecture]
Computer Design and Architecture (CDA) is the study of the structure, behavior, and implementation of computer systems. It focuses on the \textbf{Hardware/Software Interface}, bridging the gap between high-level programming languages and the physical hardware that executes them.
\end{definition}

\subsection{Core Focus Areas}

According to the UCF CDA course introduction, the primary areas of focus include:

\begin{itemize}
    \index{System Architecture}
    \item \textbf{System Architecture:} The high-level structure and organization of a computer system.
    \index{Processor Design}
    \item \textbf{Processor Design (Microarchitecture):} The internal organization of a processor, defining how it implements an Instruction Set Architecture (ISA).
    \index{Logic Design}
    \item \textbf{Logic Design:} The design of circuits using logic gates to perform operations.
    \index{Assembly Language}
    \item \textbf{Assembly Language:} Low-level programming that provides a human-readable representation of machine instructions.
    \index{Machine Instructions}
    \item \textbf{Machine Instructions:} The binary code (0s and 1s) that the hardware directly understands.
\end{itemize}

\subsection{The Levels of Abstraction}

Modern computer systems are designed using multiple levels of abstraction to manage complexity:

\begin{enumerate}
    \item \textbf{Application Software:} Written in high-level languages such as C, Java, or Python.
    \item \textbf{Systems Software:}

    \begin{itemize}
        \item \textbf{Compiler:} Translates high-level code into assembly language.
        \item \textbf{Operating System (OS):} Manages hardware resources and provides services to applications.
    \end{itemize}

    \item \textbf{Hardware:} The physical components (transistors, circuits) that execute instructions.
\end{enumerate}

\section{Basic Computer Organization}

A computer system consists of five classic components that work together to process data:

\begin{vocabulary}[Input]
Devices that feed data into the computer (e.g., keyboard, mouse, touch screen).
\end{vocabulary}

\begin{vocabulary}[Output]
Devices that present data to the user (e.g., monitor, speakers, printer).
\end{vocabulary}

\begin{vocabulary}[Memory]
Stores data and instructions currently in use.

\begin{itemize}
    \item \textbf{Volatile Memory (RAM):} Fast, but loses data when power is off.
    \item \textbf{Non-volatile Memory (SSD/HDD):} Slower, but retains data.
\end{itemize}

\end{vocabulary}

\begin{vocabulary}[Datapath]
The ``brawn'' of the processor; performs arithmetic and logical operations.
\end{vocabulary}

\begin{vocabulary}[Control]
The ``brain'' of the processor; commands the datapath, memory, and I/O devices according to the program instructions.
\end{vocabulary}

\begin{remark}
The \textbf{Central Processing Unit (CPU)} or \textbf{Processor} is the combination of the \textbf{Datapath} and \textbf{Control} units.
\end{remark}

\section{The Eight Great Ideas in Computer Architecture}

These foundational principles guide modern computer design:

\begin{itemize}
    \item \textbf{Design for Moore's Law:} Architects must anticipate where technology will be when the design finishes, rather than designing for today's constraints.
    \item \textbf{Use Abstraction to Simplify Design:} Hiding low-level details allows designers to manage complexity and improve productivity.
    \item \textbf{Make the Common Case Fast:} Focus on optimizing the most frequent operations, which yields the greatest overall performance improvement.
    \item \textbf{Performance via Parallelism:} Executing multiple tasks simultaneously (e.g., multi-core processors).
    \item \textbf{Performance via Pipelining:} Overlapping the execution of multiple instructions (like an assembly line).
    \item \textbf{Performance via Prediction:} Guessing the outcome of an operation to start work early, rather than waiting (e.g., branch prediction).
    \item \textbf{Hierarchy of Memories:} Using a combination of small, fast, expensive memory (caches) and large, slow, cheap memory (DRAM/Disk) to provide speed and capacity.
    \item \textbf{Dependability via Redundancy:} Including extra components to detect and handle failures, making the system more reliable.
\end{itemize}

\section{Moore's Law and Its Implications}

\begin{definition}[Moore's Law]
Moore's Law (proposed by Gordon Moore in 1965) predicted that the number of transistors per chip would double approximately every two years.
\end{definition}

\subsection{Key Implications}

\begin{itemize}
    \item \textbf{Exponential Growth:} Led to rapid improvements in performance and reductions in cost for decades.
    \item \textbf{The ``Power Wall'':} As clock rates increased, power dissipation became a limiting factor, ending the era of rapid single-core frequency increases.
    \item \textbf{The Current State:} Moore's Law is slowing down. Transistors are becoming harder and more expensive to shrink, leading to a shift toward \textbf{Domain-Specific Architectures (DSA)} and \textbf{Parallelism}.
\end{itemize}

\section{Performance Metrics}

Understanding performance requires distinguishing between response time and throughput.

\begin{vocabulary}[Response Time]
Also known as \textbf{Execution Time}, this is the total time required for a computer to complete a single task (user-centric).
\end{vocabulary}

\begin{vocabulary}[Throughput]
Also known as \textbf{Bandwidth}, this is the total amount of work done in a given time (system-centric).
\end{vocabulary}

\subsection{Relative Performance}

\begin{proposition}[Performance Ratio]
If Computer A is faster than Computer B, the relative performance is calculated as:
\begin{equation}
\text{Performance Ratio} = \frac{\text{Performance}_A}{\text{Performance}_B} = \frac{\text{Execution Time}_B}{\text{Execution Time}_A}
\end{equation}
\end{proposition}

\subsection{The CPU Time Equation}

\begin{theorem}[CPU Time Equation]
The fundamental formula for calculating the CPU time of a program is:
\begin{equation}
\text{CPU Time} = \text{Instruction Count} \times \text{CPI} \times \text{Clock Cycle Time}
\end{equation}
Alternatively, using the clock rate:
\begin{equation}
\text{CPU Time} = \frac{\text{Instruction Count} \times \text{CPI}}{\text{Clock Rate}}
\end{equation}
\end{theorem}

\noindent The components of this equation are determined by various factors:

\begin{itemize}
    \item \textbf{Instruction Count:} Number of instructions executed (determined by the algorithm, compiler, and ISA).
    \item \textbf{CPI (Cycles Per Instruction):} The average number of clock cycles required per instruction (determined by the microarchitecture).
    \item \textbf{Clock Cycle Time:} The duration of one clock pulse (e.g., 0.25 ns for a 4 GHz clock).
    \item \textbf{Clock Rate:} The frequency of the clock (e.g., 4 GHz).
\end{itemize}

\section{The Power Wall and Multicore Transition}

For years, performance was improved by increasing clock frequency. However, power consumption increased proportionally to $\text{Voltage}^2 \times \text{Frequency}$.

\begin{proposition}[Power Formula]
The power consumption in CMOS technology is roughly:
\begin{equation}
\text{Power} = \frac{1}{2} \times \text{Capacitive Load} \times \text{Voltage}^2 \times \text{Frequency}
\end{equation}
\end{proposition}

\begin{remark}
Due to cooling limits (the \textbf{``Power Wall''}), architects could no longer increase frequency. This forced the industry to switch from \textbf{Uniprocessors} to \textbf{Multiprocessor} (Multicore) chips, where performance is gained through parallelism rather than raw speed.
\end{remark}

\section{Classes of Computing Applications}

Modern computing is categorized into several distinct classes:

\begin{itemize}
    \item \textbf{Personal Computers (PCs):} Focus on single-user performance and low cost.
    \item \textbf{Servers:} Focus on high capacity, reliability, and security; accessed over a network.
    \item \textbf{Supercomputers:} High-end scientific/engineering calculations; hundreds of thousands of processors.
    \item \textbf{Embedded Computers:} Computers hidden inside other devices (e.g., cars, appliances), often with strict power/cost constraints.
    \item \textbf{Personal Mobile Devices (PMDs):} Battery-powered, wireless devices like smartphones and tablets.
    \item \textbf{Cloud Computing / Warehouse Scale Computers (WSC):} Giant datacenters (like Google or Amazon) providing Software as a Service (SaaS).
\end{itemize}
